\section{实验设计}

\subsection{ALU}\label{sub:alu}

\subsubsection{功能描述}
ALU模块实现32位数据的算术和逻辑运算。模块接收两个32位操作数A和B，以及3位控制码F[2:0]，输出32位运算结果Result。支持6种运算：无符号加法、减法、按位与、按位或、按位取反和小于置位（SLT）。

根据指导书，ALU有两种实现方式：
\begin{itemize}
    \item \textbf{assign实现}：使用连续赋值语句和三目运算符，适合纯组合逻辑的简单运算
    \item \textbf{always实现}：使用always@(*)块和case语句，结构更清晰，便于扩展\end{itemize}

本实验采用always实现方式，扩展支持SLT运算。

\subsubsection{接口定义}

\begin{table}[htp]
\caption{ALU模块接口定义}\label{tab:alu_signal}
\begin{center}
    \begin{tabular}{|l|l|l|p{6cm}|}
    \hline
    \textbf{信号名} & \textbf{方向} & \textbf{位宽} & \textbf{功能描述}\\ \hline \hline
    A      & Input  & 32-bit & 操作数A \\
    B      & Input  & 32-bit & 操作数B \\
    F      & Input  & 3-bit  & 运算控制码，选择ALU执行的运算类型 \\
    Result & Output & 32-bit & 运算结果，输出至七段数码管显示 \\
    \hline
    \end{tabular}
\end{center}
\end{table}

\subsubsection{逻辑控制}
ALU采用纯组合逻辑实现，通过case语句根据控制码F[2:0]选择对应运算：

\begin{table}[htp]
    \centering
    \begin{tabular}{ccc}
        \hline
        F[2:0] & 功能 & 说明 \\
        \hline
        000 & A + B (Unsigned) & 无符号加法 \\
        001 & A - B & 减法 \\
        010 & A AND B & 按位与 \\
        011 & A OR B & 按位或 \\
        100 & $\overline{A}$ & 按位取反 \\
        101 & SLT & 有符号比较，A < B 时输出1，否则输出0 \\
        110 & 未使用 & 输出0 \\
        111 & 未使用 & 输出0 \\
        \hline
    \end{tabular}
    \caption{ALU运算控制码及功能}
    \label{tab:alu_opcode}
\end{table}

SLT运算使用Verilog内置函数\texttt{\$signed()}将操作数转换为有符号数后进行比较，按照补码规则判断大小。

\subsection{4级流水线32bit全加器}\label{sub:ctl}

\subsubsection{功能描述}
实现4级流水线32位全加器。将32位加法运算拆分为4个流水阶段，每级处理8位数据的加法及进位传递，采用移位寄存器方案对齐各级的部分和输出。支持流水线暂停和刷新控制，用于模拟流水线冲突和异常处理。

流水线各级功能划分：
\begin{itemize}
    \item \textbf{第1级}：计算低8位（bit[7:0]）的加法，产生部分和sum\_s1及进位cout\_s1，剩余高位数据存入rem\_a1/rem\_b1
    \item \textbf{第2级}：计算次低8位（bit[15:8]）的加法，接收第1级进位
    \item \textbf{第3级}：计算次高8位（bit[23:16]）的加法，接收第2级进位
    \item \textbf{第4级}：计算最高8位（bit[31:24]）的加法，接收第3级进位，输出最终进位cout
\end{itemize}

由于各级部分和在不同时钟周期产生，需通过移位寄存器将各级结果对齐至同一时刻输出：
\begin{itemize}
    \item 第1级的部分和sum\_s1需延迟3个周期（经过sum1\_d2→sum1\_d3→sum1\_d4）
    \item 第2级的部分和sum\_s2需延迟2个周期（经过sum2\_d3→sum2\_d4）
    \item 第3级的部分和sum\_s3需延迟1个周期（经过sum3\_d4）
    \item 第4级的部分和sum\_s4直接作为最高8位输出
\end{itemize}

\subsubsection{接口定义}

\begin{table}[htp]
\caption{流水线全加器接口定义}\label{tab:pipe_signal}
\begin{center}
    \begin{tabular}{|l|l|l|p{5cm}|}
    \hline
    \textbf{信号名} & \textbf{方向} & \textbf{位宽} & \textbf{功能描述}\\ \hline \hline
    clk    & Input  & 1-bit  & 时钟信号 \\
    cin\_a & Input  & 32-bit & 操作数A \\
    cin\_b & Input  & 32-bit & 操作数B \\
    cin    & Input  & 1-bit  & 进位输入 \\
    stall  & Input  & 1-bit  & 流水线暂停信号，高有效 \\
    flush  & Input  & 1-bit  & 流水线刷新信号，高有效，优先级高于stall \\
    cout   & Output & 1-bit  & 进位输出 \\
    sum    & Output & 32-bit & 加法结果和 \\
    \hline
    \end{tabular}
\end{center}
\end{table}

\subsubsection{逻辑控制}

\textbf{流水线控制机制：}

\begin{enumerate}
    \item \textbf{正常流动}：stall=0且flush=0时，每个时钟周期数据向下一级推进。各级并行计算各自负责的8位加法，进位在各级间传递
    \item \textbf{暂停}：stall=1时，所有流水线寄存器保持当前值不变，各级计算暂停
    \item \textbf{刷新}：flush=1时，所有流水线寄存器清零，包括各级部分和、进位、剩余数据以及移位寄存器。stall和flush同时有效时，flush优先
\end{enumerate}

\textbf{仿真控制时序：}
\begin{itemize}
    \item 第10个时钟周期后，stall信号有效2个周期，模拟第2级流水线暂停，观察流水线各级寄存器保持情况
    \item 第15个时钟周期时，flush信号有效1个周期，模拟第3级流水线刷新，观察流水线寄存器清零和重新流动
\end{itemize}

\textbf{进位传递与结果拼接：}

每级加法使用\{进位, 部分和\}的拼接格式：\{cout\_sN, sum\_sN\} <= \{1'b0, a\_part\} + \{1'b0, b\_part\} + carry\_in。进位cout\_s1从第1级传递至第2级，cout\_s2传递至第3级，cout\_s3传递至第4级。最终进位cout由第4级产生。由于各级部分和在流水线中处于不同阶段，通过移位寄存器对齐后拼接为32位最终结果：sum[7:0]来自延迟3拍的sum\_s1，sum[15:8]来自延迟2拍的sum\_s2，sum[23:16]来自延迟1拍的sum\_s3，sum[31:24]直接来自sum\_s4。
